% =========================================================================
% Prose rewritten by Aryan Padarthi, 4 October 2026.
% Mathematics (displayed formulas, theorem statements) restored from the
% main paper. Nine corrections marked % [FIX n] -- see the message that
% accompanied this file.
% =========================================================================
\documentclass[11pt]{article}
\usepackage[margin=1in]{geometry}
\usepackage{amsmath,amssymb,amsthm,hyperref}

\newtheorem{theorem}{Theorem}
\newtheorem{lemma}[theorem]{Lemma}
\newtheorem{corollary}[theorem]{Corollary}
\theoremstyle{definition}
\newtheorem{remark}[theorem]{Remark}

\title{The zero forcing number of $P(n,3)$}
\author{Aryan Padarthi}
\date{October 2026}

\begin{document}
\maketitle

\begin{abstract}
In Theorem~3.6, Rashidi, Shajareh Poursalavati and Tavakkoli write that
$Z(P(n,3))=8$ for all $n$ at least $12$. Krishnan found that this statement is
false for $n=12$, where $Z(P(12,3))=7$, calculated $Z(P(n,3))$ exhaustively for
$n$ from $7$ to $20$, and conjectured that $Z(P(n,3))=8$ for every $n$ at least
$13$. He remarked that what was missing was a proof of a lower bound which holds
for all large $n$. Here we provide that proof, and show that the conjecture
holds.
% [FIX 1] "which happens precisely if T is the identity" attached to "at most
% 8"; the condition is for EQUALITY, not for the inequality.
Our proof is a reduction. We observe that removing the interior coordinates of
any matrix with the pattern of $P(n,3)$ produces a single linear recurrence of
order $8$, so the nullity equals the dimension of the fixed space of the
monodromy $T$, and is at most $8$, with equality precisely when $T$ is the
identity.
% [FIX 2] tiles do not "produce n"; they give the conclusion for every n.
We then construct matrices achieving $8$ by concatenating tiles, each having a
fixed pattern for docking, so that their transfer products depend only on their
interiors. There are tiles of every length from $21$ up to but not including
$42$, and together they settle every $n$ at least $21$ at once. We certify each
tile by evaluating a Krawczyk contraction in exact rational arithmetic. With
per-$n$ certificates for $n$ from $17$ to $20$ and an exhaustive search for $n$
from $13$ to $16$, we have proved $Z(P(n,3))=8$ for every $n$ at least $13$.
Here $13$ is optimal, because $Z(P(11,3))=Z(P(12,3))=7$.
\end{abstract}

\section{Introduction}

Let $G$ be a graph whose vertices are each coloured black or white. The colour
change rule is: if a black vertex has exactly one white neighbour, that
neighbour becomes black. Let $S$ be a set of vertices. Then $S$ is zero forcing
if colouring $S$ black and then repeatedly applying the rule until it stalls
colours all vertices. Let $Z(G)$ denote the smallest possible size of a zero
forcing set. If $M(G)$ denotes the maximum nullity among all symmetric real
matrices having exactly the edge set of $G$ as their nonzero off-diagonal
entries, then $M(G)$ is at most $Z(G)$ \cite{aim}.

The generalized Petersen graph $P(n,k)$, for $n\ge3$ and $1\le k<n/2$, has
vertices $u_0,\dots,u_{n-1}$ and $v_0,\dots,v_{n-1}$, where $u_i$ is joined to
$u_{i+1}$, $v_i$ is joined to $v_{i+k}$, and $u_i$ is joined to $v_i$. All
indices are taken modulo $n$. $P(n,k)$ is cubic.

In their Theorem~3.6, Rashidi, Shajareh Poursalavati and Tavakkoli \cite{rst}
write that $Z(P(n,3))=8$ for any $n\ge12$. Krishnan \cite{k} observed that this
claim is false for $n=12$, and gives an explicit set of seven elements which
constitutes a zero forcing set for $P(12,3)$. He also points out that an error
was made in the proof of the result in the published paper, since it only rules
out six-element sets in its case analysis.
% [FIX 3] "does not change after n = 12" reads as though Z(12)=8; the values
% stabilise FROM 13, and Z(12)=7.
By brute force he determines $Z(P(n,3))$ for $n=7,\dots,20$, finds that the
values stabilise at $8$ from $n=13$ onward, and thus concludes:

\medskip
\noindent\textbf{Conjecture 5.} $Z(P(n,3))=8$ for any $n\ge13$.
\medskip

The obstacle is a matching lower bound valid for all $n$: the upper bound is
available for every $n$ at least $9$, and only finitely many lower bounds can
come from computation.

\begin{theorem}\label{thm:main}
$Z(P(n,3))=8$ for every $n\ge13$, and $13$ is optimal, since
$Z(P(11,3))=Z(P(12,3))=7$. Moreover $Z(P(n,3))=M(P(n,3))=8$ for every $n\ge17$.
\end{theorem}

This is Conjecture~5 of Krishnan.

The upper bound $Z(P(n,3))\le8$ for $n$ at least $9$ is a rotation argument and
is not new. The content is the lower bound, and all of it is in
Sections~\ref{sec:red}--\ref{sec:cert}. Section~\ref{sec:red} reduces the kernel
of any matrix carrying the pattern of $P(n,3)$ to the periodic solutions of a
linear recurrence of order $8$, which bounds the nullity by $8$ and identifies
exactly when the bound is attained: when the monodromy of the recurrence is the
identity. Section~\ref{sec:tiles} makes that condition achievable for every
large $n$ simultaneously, by assembling the weight sequence out of tiles whose
transfer products multiply independently. Section~\ref{sec:cert} certifies the
tiles. Section~\ref{sec:proof} assembles the proof.

\section{Setup}\label{sec:setup}

Fix $n$ and $k$ with $n\ge2k+3$. Let $A$ be a real matrix indexed by the
vertices of $P(n,k)$, with a nonzero entry in position $xy$ if and only if $xy$
is an edge, for $x\ne y$; the diagonal is unrestricted. We do not assume $A$ is
symmetric. Write
\[
  b_i=A_{u_iu_{i+1}},\quad b'_i=A_{u_{i+1}u_i},\quad
  c_i=A_{u_iv_i},\quad c'_i=A_{v_iu_i},\quad
  e_i=A_{v_iv_{i+k}},\quad e'_i=A_{v_{i+k}v_i},
\]
with diagonal entries $a_i$ at $u_i$ and $d_i$ at $v_i$. All of
$b,b',c,c',e,e'$ are nonzero.

We do not need symmetry to prove that the nullity of $A$ is at most $Z(G)$. We
only need to know that there is a nonzero entry on every edge in both
directions. If we have $Ax=0$ and $x$ is zero on a set $B$, and there is a
vertex $u$ in $B$ such that $u$ has exactly one neighbour $w$ not in $B$, then
all the terms in the row of $Ax=0$ at $u$ will be zero except the one involving
$x_w$, which has a nonzero coefficient, so $x_w=0$.
% [FIX 4] "gives an injection onto the kernel of A" -- wrong direction and
% "onto" would be surjectivity. The restriction map is injective ON the kernel.
We can show by induction over a forcing process that for any zero forcing set
$S$, the restriction map $x\mapsto x|_S$ is injective on the kernel of $A$. So
\begin{equation}\label{eq:mz}
  \operatorname{null}A\ \le\ Z(G).
\end{equation}

\section{The reduction}\label{sec:red}

This allows us to combine the outer and inner equations of the kernel of $A$
into one recurrence along the outer cycle. Taking $k=3$, we find that the
recurrence is of order $8$, which is the bound we want to match.

\begin{theorem}[Reduction]\label{thm:red}
Let $n\ge2k+3$ and for $x\in\mathbb R^{\mathbb Z_n}$ define
\[
  (Lx)_i=-\frac{b'_{i-1}x_{i-1}+a_ix_i+b_ix_{i+1}}{c_i},
\]
obtained by solving the outer equation for the inner coordinate. Then the map
taking $x$ to the pair $(x,Lx)$ is a linear isomorphism from
\[
  \Big\{x:\ \textstyle\sum_{p=-k-1}^{k+1}\gamma_{i,p}\,x_{i+p}=0
  \ \text{ for all } i\in\mathbb Z_n\Big\}
\]
to the kernel of $A$, where $\gamma_{i,p}\ne0$ only for
$p\in\{-k-1,-k,-k+1\}\cup\{-1,0,1\}\cup\{k-1,k,k+1\}$, and the two extreme
coefficients are
\[
  \gamma_{i,k+1}=-\frac{e_i\,b_{i+k}}{c_{i+k}}\ne0,
  \qquad
  \gamma_{i,-k-1}=-\frac{e'_{i-k}\,b'_{i-k-1}}{c_{i-k}}\ne0 .
\]
Therefore, with $T$ the monodromy of the recurrence on the state
$(x_{i-k-1},\dots,x_{i+k})$, the nullity of $A$ is the dimension of the fixed
space of $T$, is at most $2k+2$, and equals $2k+2$ if and only if $T$ is the
identity.
\end{theorem}

\begin{proof}
Every spoke weight is nonzero, so the equation of the kernel at $u_i$ has a
unique solution for the inner coordinate at $i$, a three-term expression in $x$.
Putting this into the equation at $v_i$, namely
$e'_{i-k}y_{i-k}+d_iy_i+e_iy_{i+k}+c'_ix_i=0$, gives the formula: the three
substituted terms give the three index triples $\{i-k-1,i-k,i-k+1\}$,
$\{i-1,i,i+1\}$ and $\{i+k-1,i+k,i+k+1\}$, and the two extreme coefficients are
as claimed, both nonzero since the outer, spoke and inner weights are. Every
vector $x$ solving all these equations gives a vector in the kernel.

Both extreme coefficients are nonzero, so the relation is a recurrence of order
exactly $2k+2$, which is solvable forward and backward, with solutions on the
integers giving a space of dimension $2k+2$, which the monodromy acts on, and
whose $n$-periodic solutions are exactly the fixed vectors of the monodromy.
This proves that the nullity of $A$ is the dimension of that fixed space. This
is at most $2k+2$ and equals $2k+2$ if and only if the monodromy is the
identity.
\end{proof}

% [FIX 5] "the monodromy's order could be less than 2k+2" -- the monodromy is a
% matrix; it is the RECURRENCE whose order can drop. Also merged the two
% sentences, which said the same thing twice.
This proof requires that $n$ be at least $2k+3$. Otherwise the index sets above
collide modulo $n$, the extreme coefficients are no longer given by the explicit
formulas, and for $n$ equal to $2k+1$ or $2k+2$ the recurrence can have order
less than $2k+2$.

\begin{corollary}\label{cor:ceiling}
$M(P(n,3))\le8$ for $n\ge9$.
\end{corollary}

For $k=3$ the order of the recurrence is $8$, and the size of the forcing set
produced by the rotation argument giving the upper bound is also $8$. The upper
bound and the ceiling of the matrix method are identical, with no slack: there
exists a matrix on the pattern of $P(n,3)$ with nullity $8$ if and only if the
monodromy of its reduced recurrence is the identity, and such a matrix proves
the lower bound.

\section{Tiles}\label{sec:tiles}

The docking pattern consists of all weights equal to $1$ in the outer, spoke and
inner positions, and zero in the two diagonal positions. A tile of length $\ell$
is a group of $\ell$ consecutive positions in which the first $k+1$ and last
$k+1$ positions have the docking pattern, and the other positions make up its
interior. The tile product is the product of the transfer matrices over the
positions of the tile.

\begin{lemma}[Tile independence]\label{lem:tile}
If the weight sequence is a concatenation of tiles sharing one docking pattern,
each tile product depends only on that tile's interior weights, and the
monodromy is the product of the tile products.
\end{lemma}

\begin{proof}
By the reduction theorem, each transfer matrix depends on the coefficients of
one reduced equation, involving weights only at indices in a window of width
$2k+2$ around its position. If the tile is in a block and the position is in
that block, then the window will be in the block padded out by $k+1$ on each
side. The positions below the block will be in the last $k+1$ positions of the
previous tile and the positions above the block will be in the first $k+1$
positions of the next tile; these also have the docking pattern, and so does the
tile itself in its first and last $k+1$ positions. Everything the transfer
matrix sees is either docking or interior to this tile. The second statement
says that the monodromy product associates when regrouped over the tiles.
\end{proof}

% [FIX 6] The proof of the tiling theorem appeared BEFORE its statement in the
% rewrite. Reordered.
\begin{theorem}[Tiling]\label{thm:tiling}
Fix $k$ and $L\ge2(k+1)$. Suppose that for every length from $L$ up to but not
including $2L$ there are interior weights, all required entries nonzero, such
that the tile product equals the identity. Then
\[
  Z(P(n,k))=M(P(n,k))=2k+2 \qquad\text{for every } n\ge L .
\]
\end{theorem}

\begin{proof}
Each $n$ at least $L$ is a sum of integers from that range: if $n$ is below $2L$
take the single term $n$, otherwise take a term $L$ and recurse on $n-L$. We can
concatenate together the corresponding tiles, giving a weight sequence with
monodromy that by Lemma~\ref{lem:tile} is the product of the tile products, and
hence the identity. All weights are nonzero, so it represents the pattern of
$P(n,k)$, and from the reduction theorem the nullity is $2k+2$.
% [FIX 7] "This can be seen using the inequality from Section 2 and the rotation
% bound. Both are 2k+2." -- that is a separate deduction, not a restatement.
Hence $M(P(n,k))\ge2k+2$; combining this with \eqref{eq:mz} and the rotation
bound $Z\le2k+2$ gives $M(P(n,k))=Z(P(n,k))=2k+2$.
\end{proof}

% [FIX 8] "the first possible lengths here to solve is several hundred" --
% garbled; it is the uncovered range that is several hundred.
We have proven this without gaps because we used a full interval. If the lengths
of two tiles were coprime, then they could cover only nonnegative integer
combinations of the two, and every $n$ below a certain product would not be
covered; for the smallest solvable lengths here, that uncovered range is already
several hundred. Our hypothesis is a finite set of independent finite problems.

\section{Certification}\label{sec:cert}

% [FIX 9] "As an example consider k = 3" -- k = 3 is the case of the paper, not
% an illustration.
For $k=3$ we take $L=21$ and exhibit identity tiles of all lengths from $21$ to
$43$ inclusive, which includes the needed range. That is $23$ tiles.

To certify directly that the tile product is the identity would be to do
interval arithmetic on a rational map. We do not need to do that. Applying the
reduction theorem to the cyclic matrix formed from one tile tells us that the
monodromy of that matrix is the tile product. Therefore the tile product is the
identity if and only if the nullity of that matrix is $2k+2$, and that last
statement is about a bilinear system. Say $K$ is a matrix with columns spanning
a possible kernel, written in graph form. The property is that $A(w)K(X)=0$, and
its Jacobian is affine in the unknowns.

% [FIX 10 -- the one that matters] The rewrite said removing the dependencies
% leaves "a system which is still a consequence of the original system". That
% inverts the point: the whole reason for choosing the subsystem structurally
% is that it is EQUIVALENT to AK = 0, not merely implied by it. As written it
% would have weakened the central rigour claim of the note.
We take the square subsystem not because it is the maximum-rank selection one
could make, but because it is chosen structurally. The product
$K^{\!\top}\!AK$ is symmetric, and its antisymmetric part accounts for precisely
$\binom{8}{2}=28$ dependencies among the equations, so removing those leaves a
system equivalent to $AK=0$, rather than a weaker consequence of it.

For each tile we then perform a Krawczyk contraction to certify it. With $Y$ an
approximate inverse of the square Jacobian $J$ at a numerically located point,
$L_{\!J}$ a Lipschitz bound for $J$ in the $\infty$-norm, and $\delta$ a box
radius,
\[
  \alpha=\|I-YJ\|_\infty+\|Y\|_\infty L_{\!J}\,\delta<1
  \qquad\text{and}\qquad
  \|Yf\|_\infty+\alpha\delta\le\delta
\]
together establish that a true real solution exists in that box.

Every number appearing here is evaluated exactly: the box centre is dyadic; the
system is bilinear with integer coefficients, so the value and Jacobian at the
centre are both exactly rational; and the approximate inverse is arbitrary, so
rounding it to a rational costs nothing. Across all $23$ tiles the contraction
constant is at most $1.0\times10^{-9}$, and the nullity of each certified matrix
is exactly $8$. In addition, every docking weight is equal to its prescribed
integer. Finally, none of the above uses any nullity calculation with
floating-point numbers.

Performing exactly the same process per $n$ rather than per tile gives us a
proof that $Z$ and $M$ are both equal to $8$ for every $n$ between $17$ and
$33$, that is, for seventeen consecutive numbers including the primes
$17,19,23,29$ and $31$.

\section{Proof of the theorem}\label{sec:proof}

\begin{proof}[Proof of Theorem~\ref{thm:main}]
First, the upper bound is correct for all $n\ge9$, by the rotation argument.

For $n\ge21$ we have $Z=M=8$, since there are certified identity tiles of all
lengths in the right range given by Section~\ref{sec:cert}, so
Theorem~\ref{thm:tiling} applies with $L=21$.

For $n$ from $17$ to $20$: these are covered by the per-$n$ certificates of
Section~\ref{sec:cert}, which give $Z=M=8$ on the range $17$ to $33$.

Finally, for $n$ in $[13,16]$, exhaustive search on subsets of vertices shows
there is no zero forcing set of size $7$. Hence $Z=8$ by the upper bound. We
deduce $Z$ only, and do not assert that $M=8$ here.

Optimality of $13$: exhaustive search gives $Z(P(11,3))=Z(P(12,3))=7$.
\end{proof}

\begin{remark}
The sequence of values for $n$ from $7$ to $14$ is $6,6,6,8,7,7,8,8$. So $Z=8$
first occurs at $n=10$, well below the threshold $13$ beyond which it always
holds. These are two different quantities.
\end{remark}

\begin{remark}
By the reduction theorem the lower bound is attained precisely when a monodromy
equals the identity, a condition with no arithmetic in it. That is why the proof
reaches all large $n$, including primes, and not a congruence class.
\end{remark}

\section*{Independence and related work}

This work was carried out independently. The lower bound for $n\ge17$ was
obtained on 3 September 2026 and the determination of $Z(P(n,3))$ for every $n$
on 22 September 2026; both dates are recorded in the version history of the
project repository. I re-ran the verification suite for every computational
claim in this note before posting.

After obtaining these results I learned of \cite{lean}, a pull request to the
repository \texttt{the-omega-institute/trureturing} (\#10129, opened and merged
on 26 September 2026), which states and claims a Lean proof of the same
conjecture. I looked over it myself. What I verified is that it defines the
statement
\[
  \texttt{claim} \;:\; \forall\, n,\ 13\le n \implies
  \texttt{zeroForcingNumber (gp n 3)} = 8,
\]
which is Conjecture~5, that it contains a declaration \texttt{theorem result :
claim}, and that its Lean sources contain no \texttt{sorry} and no additional
\texttt{axiom} declarations. I did not build it, so I make no claim about
whether it compiles, and I make no other assessment of it.

\begin{thebibliography}{9}
\bibitem{aim} AIM Minimum Rank -- Special Graphs Work Group, Zero forcing sets
  and the minimum rank of graphs, \emph{Linear Algebra Appl.} \textbf{428}
  (2008), 1628--1648.
\bibitem{k} A. Krishnan, A correction to the zero forcing number of the
  generalized Petersen graphs $P(n,3)$, arXiv:2607.19412 (2026).
\bibitem{lean} \texttt{the-omega-institute/trureturing}, pull request \#10129,
  ``Prove Krishnan's Conjecture 5: the zero forcing number of $P(n,3)$ is 8 for
  every $n\ge13$'', merged 26 September 2026.
  \url{https://github.com/the-omega-institute/trureturing/pull/10129}
\bibitem{rst} S. Rashidi, N. Shajareh Poursalavati, M. Tavakkoli, Computing the
  zero forcing number for generalized Petersen graphs, \emph{J. Algebra Combin.
  Discrete Struct. Appl.} \textbf{7}(2) (2020), 183--193.
\end{thebibliography}

\end{document}
